% !TeX root = ../main.tex
\chapter{地震应急资源服务网络建模与统一分析框架}
\label{chap:framework}

\section{应急资源服务网络的系统边界}
\subsection{道路、资源、需求与服务之间的关联}
\subsection{物资交付服务的基本分析单元}
\subsection{物理状态、观测信息与决策行动的区分}

\ThesisBlank

\section{研究区域与数据基础}
\subsection{泸定县及其区域互援联系}
\subsection{道路网络与灾害情景数据}
\subsection{人口需求与应急资源数据}

\ThesisBlank

\section{供需空间表征与参数标定}
\subsection{区域需求与资源能力的相对匹配}
\subsection{相对指标与服务模型参数的衔接}
\subsection{参数不确定性与情景分层}

\ThesisBlank

\section{道路与资源约束下的服务过程模型}
\subsection{道路通行时间、容量与路径约束}
\subsection{资源库存、在途流量与需求满足}
\subsection{需求时限与服务状态更新}

\ThesisBlank

\section{服务损失与动态韧性的统一测度}
\subsection{按时满足程度与服务缺口}
\subsection{累积服务损失与韧性轨迹}
\subsection{效率、公平与行动成本的评价边界}

\ThesisBlank

\section{理论分析与实验验证的统一设计}
\subsection{受控网络、泸定路网与灾损情景}
\subsection{共同基准、资源预算与信息条件}
\subsection{独立重复、配对比较与不确定性报告}
\subsection{机制归因、失效情景与外推边界}

\ThesisBlank

\section{本章小结}
\ThesisBlank
